% ============================================================================
% Fichier cible : ARIEL_LOGMINER_MEMOIRE_FINAL/main.tex, chapitre étoilé « Abstract »
% Lignes à remplacer : 399 à 499 (de « This thesis presents Ariel \logminer{}, a
%   lightweight multi-agent architecture » jusqu'à « rather than a new
%   machine-learning algorithm. »). Conserver les lignes 396-397 et le bloc
%   des mots-clés (lignes 501-507).
% Auteur : Agent 1 (intégrité factuelle), chantier 2026-11-05.
% Statut : provisoire ; traduction fidèle de main.tex__resume_fr.tex (même
%   source pour chaque chiffre). Paragraphe « hypotheses » modifiable par
%   l'Agent 2 ; paragraphes architecture reformulables par l'Agent 5.
% ============================================================================

This thesis presents Ariel \logminer{}, a lightweight multi-agent architecture
for heterogeneous system and network log analysis, whose distribution is
evaluated in a laboratory topology. An agent is defined as a software entity
with declared capabilities, state, memory, and a local policy. When assigned
a bounded objective, it can evaluate available tasks, propose or refuse their
execution, invoke compatible handlers, chain subtasks, resume persistent state,
and delegate an operation requiring another capability. Negotiation follows an
explicit Contract Net Protocol over a local bus or Redis Streams. Allocation
and distribution are evaluated through laboratory campaigns, while bounded
goal pursuit is functionally validated. General intelligence, autonomous goal
definition, and operation as an industrial SOC platform remain outside the
scope of this work.

The evaluation particularly highlights the influence of the experimental
protocol on measured performance. On CICIDS2017, a Random Forest evaluated
with a stratified random split reaches a mean F1-score of $0.995142$ over five
seeds. When each of the five attack scenarios is held out from training, the
mean F1-score falls to $0.157744$. Among five lightweight models compared
under this holdout, Logistic Regression achieves the highest mean F1-score,
$0.233670$, while HistGradientBoosting reaches a macro PR-AUC of $0.569002$
with a near-zero F1-score at the selected threshold. Changing the classifier
therefore alters some metrics without eliminating inter-scenario fragility.

For HDFS and BGL, the protocol strictly separates training, validation, and
test partitions. Drain3 is fitted only on the training partition and frozen
before validation and test. At the \texttt{block\_id} level, Histogram reaches
an F1-score of $0.868373$ (95\% CI $[0.856771;0.879493]$) on the
$115{,}013$ HDFS test blocks, including $1{,}679$ positives, with a recall of
$1$ and a precision of $0.767367$, identical across five sampling seeds. On
BGL, F1 is $0.913698$, but 89.8081\% of test events carry a template unseen
in training, and the Histogram decision coincides exactly with the rule
``severity absent from normal training''. This score therefore measures
severity novelty rather than a model of log dynamics; the ``unknown
template'' rule reaches only $0.277885$.

On CSE-CIC-IDS2018, a February 15 versus February 16 split yields a mean
F1-score of $0.999212$ for Logistic Regression and $0.401375$ for Random
Forest. However, \texttt{Dst Port} alone reaches $0.999720$, a score also
obtained by a fixed destination-port rule (below 1024), and the reverse
transfer, from February 16 to February 15, does not exceed $0.234195$. The
primary score therefore rests on a shortcut specific to the test day.

The multi-format validation reads, parses, and normalizes $7{,}001/7{,}001$
units, including a single Apache fixture. This result measures functional
coverage of the tested formats; filling of the common schema ranges from seven
standard fields out of seven (syslog, BGL) to none (tabular network flows). An
open-set policy calibrated only on known sources and then frozen rejects
$20/20$ unknown files and wrongly rejects $3/46$ known files; all 43 accepted
known files are correctly routed, and the decision margin is not a calibrated
probability. No comparison with a pipeline without routing was run: the
predictive benefit of specialized routing is not established.

The controlled architecture benchmark covers four architectures, four loads,
and ten repetitions, yielding 160 runs and $264{,}000$ treatments with no
failure. On the protocol's microtasks, the monolith dominates and the agent
architectures add an orchestration cost. Under a paired failure occurring
after persistence, the agent architectures duplicate no effect, unlike the
monolith and the centralized workers; equipping the latter with the same
idempotency store also removes duplicates.

A one-hour endurance campaign on two virtual machines, Debian and Ubuntu,
connected to a Redis instance centralized on the laboratory host, completes
$1{,}800$ out of $1{,}800$ tasks, split 451 and $1{,}349$ according to the
configured capability profiles; a controlled stop after persistence leads to
a reassignment and a replay with no duplicated effect. Under saturation, the
capacity of this topology lies between 8 and 16 tasks/s, and the second
machine processes only 1.23\% of the successful tasks: it does not increase
the observed capacity. High availability was not evaluated.

The correlation mechanism, evaluated on a controlled synthetic corpus,
obtains a precision of $0.764706$, a recall of $0.866667$, and a pairwise
F1-score of $0.812500$. On larger generated corpora, its recall falls to
$0.5973$; when 30\% of incidents chain several event categories, a variant
reduced to host and time window obtains a higher F1-score ($0.7268$ versus
$0.6496$).

In the multi-source CNP campaign, condition M actually runs compatible
artifacts on $1400/1401$ units, with one Apache fallback and no error. Its
mean latency is $0.105420$~s, compared with $0.031027$~s for the heuristic H
condition. The lack of homogeneous labels precludes a global multi-source
accuracy; the Linux/auth F1 of $0.390244$ remains descriptive of the selected
source.

The repository also contains a model-comparison procedure and a dry-run mode.
The tests cover some branches, but not a complete and repeated
promotion/rejection cycle. Its overall benefit was not evaluated.

No hypothesis is validated: three (H1, H2, and H4) are partially supported
and two (H3 and H5) remain undecided, for lack of a no-routing comparator for
the former and of a demonstrated quantitative memory gain for the latter. The
answers to the research questions remain partial; the methodological question
(RQ6) is the best covered within the laboratory scope. Ariel \logminer{}
therefore contributes primarily an architectural and methodological
integration of a local, modular, and experimentally traceable
heterogeneous-log analysis pipeline, rather than a new machine-learning
algorithm.
